\clearpage
\fancyhead[L]{\small CAGNOLI. STERNVERZEICHNISS.}
\hypertarget{native-501-lower}{}\label{native-501-lower}
\begin{center}
\small Göttingische gelehrte Anzeigen. Stück 77. Seite 761..763. 1808 Mai 14.\\[3pt]
\rule{.72\linewidth}{.3pt}
\end{center}

\emph{Catalogue de \textup{501} étoiles suivi des tables relatives d'Aberration et de Nutation, par Antoine Cagnoli.} Modena 1807. 280 Seiten in Quart.

Dieses Sternverzeichniss war schon im X. Bande der \emph{Memorie di Matematica e di Fisica della Societa Italiana} abgedruckt, und erscheint hier verbessert und vermehrt. Jener erste Abdruck war auch von historischen Nachrichten über die Entstehung des Catalogs begleitet, wovon wir hier Einiges ausheben. \textsc{Cagnoli} hatte die Verfertigung eines Sternverzeichnisses aus eignen Beobachtungen schon im Jahre 1783 an seinem damaligen Aufenthaltsorte Paris angefangen, und nachher diese Arbeit in den Jahren 1788\mbox{. . .}1792 zu Verona mit denselben Instrumenten vollendet. Der Plan, welchen er sich dabei vorgesetzt hatte, ging dahin, nördlich vom Aequator in jeder Zone von einem halben Grad Ausdehnung in der Declination wenigstens vier zuverlässige Bestimmungen von Sternen zu geben, die in der geraden Aufsteigung ungefähr 6 Stunden von einander abständen. Hieraus ist die anfangs auffallende Auswahl der Sterne zu erklären, da keineswegs die hellsten in jedem Sternbilde ausgehoben sind, sondern statt dieser sich oft Sterne bis zur sechsten Grösse in dem Verzeichnisse befinden. Ganz konnte indessen \textsc{Cagnoli} doch diesem Plane nicht treu bleiben, theils weil nicht überall schickliche Sterne zu finden waren, theils weil Gesundheitsumstände und andre Störungen ihn daran hinderten. \textsc{Cagnoli} wollte bei seinem Verzeichniss nichts von Andern entlehnen; er bestimmte also, um die absoluten geraden Aufsteigungen zu erhalten, die von $\alpha$ im Fuhrmann durch 24 unmittelbare Vergleichung mit der Sonne, und fand hieraus ein mit \textsc{Maskelyne}'s neuester Angabe bis auf $1''$ stimmendes Resultat. Die Instrumente, womit die ganze Arbeit ausgeführt ist, sind ein beweglicher Quadrant von drei Fuss Halbmesser, dessen achromatisches Fernrohr zwei Zoll Oeffnung hat; ein achromatisches 3$\frac{1}{2}$fussiges Mittagsfernrohr mit 28 Linien Oeffnung, beide von \textsc{Megnié}, und eine vortreffliche, von \textsc{Robins} in Paris verfer%
\SourceLeaf{502}%
tigte, Pendeluhr. Alle diese Instrumente befinden sich gegenwärtig auf der Sternwarte zu Mailand. Einige Sterne sind übrigens nach Beobachtungen von \textsc{Cesaris} in Mailand bestimmt, und im Catalog durch ein besonderes Zeichen unterschieden.

Von den 501 Sternen des Verzeichnisses sind die meisten nördliche, nur 28 sind südliche. Ihre Stellungen sind für 1800 angegeben, und jeder Rectascension und Declination ist die Anzahl der Beobachtungen und der Unterschied der am meisten abweichenden mit beigefügt. Die Unterschiede von \textsc{Piazzi}'s Bestimmungen gehen bei dem grössern Theile der Sterne nur auf wenige Secunden, und es ist wirklich ein Beweis von vorzüglicher Geschicklichkeit und Sorgfalt, dass \textsc{Cagnoli} mit seinen, für einen solchen Zweck doch nur mittelmässigen, Hülfsmitteln diesen Grad von Schärfe hat erreichen können. Am Ende des Verzeichnisses sind sämmtliche Sterne noch einmal nach ihren Abständen vom Nordpol geordnet, welches für diejenigen, die diese Sterne zur Vergleichung anwenden wollen, sehr bequem ist.

Vorzüglich schätzbar sind die dem neuen Abdruck des Catalogs beigefügten speciellen Aberrations- und Nutations-Tafeln für die meisten dieser Sterne; erstere sind bereits nach der neuesten Bestimmung der Constante der Aberration von \textsc{Delambre} ($20''25$) berechnet, letztere hingegen gründen sich noch auf die Annahme der halben grossen und kleinen Axe der Nutations-Ellipse zu $9''0$ und $6''7$, und alle Zahlen derselben müssen demnach um den neunten oder vierzehnten Theil vergrössert werden, wenn sie mit \textsc{Laplace}'s oder \textsc{von Zach}'s neuesten Angaben in Uebereinstimmung gebracht werden sollen.
